\documentclass[11pt]{article}
\usepackage[margin=1in]{geometry}
\usepackage{hyperref}
\usepackage{booktabs}
\usepackage{float}
\usepackage{graphicx}
\usepackage{amsmath}
\usepackage{enumitem}
\usepackage{titlesec}

\titleformat{\section}{\large\bfseries}{\thesection.}{0.5em}{}
\titleformat{\subsection}{\normalsize\bfseries}{\thesubsection.}{0.5em}{}

\title{\textbf{Kian: A Fully Offline Voice Assistant for Children's Education}}
\author{Aaron Schweiger, 2026}
\date{\vspace{-2em}}

\begin{document}
\maketitle

\begin{abstract}
Kian is an open-source, fully offline voice assistant designed for children's education~\cite{kian2026}.
Running entirely on an NVIDIA Jetson Orin, it requires no internet connection and no parent-managed accounts---children interact through speech alone.
A multi-stage content filtering pipeline reduces exposure to harmful material, while a heuristic LaTeX-to-speech post-processing system makes audio-based mathematical question answering feasible.
This paper describes the architecture, design rationale, and key tradeoffs of the system.
\end{abstract}

\section{Motivation}

Voice assistants have become a natural interface for children to explore questions, practice concepts, and satisfy curiosity.
However, existing solutions depend on cloud services, raising concerns about privacy, availability, and access control.
While recent work has explored pruning and quantizing language models for on-device deployment~\cite{liu2024comparing},
formal evaluation of complete offline voice assistant pipelines designed for children---where content safety, speech-only interaction,
and age-appropriate language are first-class requirements---remains limited.

\begin{itemize}[nosep]
  \item \textbf{Privacy.}
    Children's voice data is sensitive.
    Cloud-based systems transmit raw audio to remote servers, creating risks under regulations such as COPPA.
    An offline system ensures that no audio or conversation data ever leaves the device.
  \item \textbf{Connectivity.}
    Many educational settings---rural schools, outdoor classrooms, vehicles, areas with unreliable infrastructure---lack consistent internet access.
    A fully on-device system operates anywhere.
  \item \textbf{Autonomy.}
    Children should be able to interact with an educational tool without needing a parent's phone, a managed account, or an app store.
    Kian is controlled entirely through speech: changing the persona, adjusting volume, setting grade level, and resetting context are all voice commands.
\end{itemize}

\section{Related Work}

A growing body of HCI research has examined children's interactions with voice-based conversational agents;
Garg et al.~\cite{garg2022last} survey a decade of this work,
identifying persistent concerns around privacy, developmental impact, and the absence of child-specific design in commercial assistants.
Several open-source projects have pursued offline voice processing---notably
Open Voice OS~\cite{ovos} (continuing the discontinued Mycroft platform),
which has also been used to build offline assistants on the NVIDIA Jetson~\cite{ovos_jetson},
and Rhasspy~\cite{rhasspy}, whose author also created the Piper TTS engine~\cite{piper_tts} used in this work.
These systems were originally designed around intent parsers rather than generative LLMs,
though recent integrations have added LLM support.
The challenge of rendering mathematics as speech has been studied since Raman's AsTeR system~\cite{raman1998audio},
which first converted \LaTeX{} to structured audio;
Kian's post-processor draws on this tradition,
using explicit delimiters for disambiguation in the style of the MathSpeak convention.

\section{System Architecture}

Kian is structured as an asynchronous pipeline of five stages, all running on-device:

\begin{enumerate}[nosep]
  \item \textbf{Voice Activity Detection (VAD).}
    \texttt{Silero VAD}~\cite{silero_vad} (ONNX, ${\sim}2$\,MB) monitors the microphone stream at 48\,kHz,
    downsampled to 16\,kHz for inference.
    Speech segments are extracted with a 1.2\,s silence threshold.
    A 30-second maximum prevents unbounded memory growth from sustained noise.
  \item \textbf{Speech-to-Text (STT).}
    \texttt{Faster-whisper}~\cite{faster_whisper}, an optimized reimplementation of Whisper~\cite{radford2023whisper},
    with the \texttt{base.en} model transcribes detected speech segments using int8 quantization.
  \item \textbf{Language Model (LLM).}
    Qwen3-4B\footnote{Specifically the Qwen3-4B-Instruct-2507 variant, which uses non-thinking mode only and does not generate \texttt{<think>} blocks.}~\cite{yang2025qwen3},
    quantized to Q4\_K\_M (4-bit), runs via llama.cpp~\cite{llamacpp_python} with full GPU offload.
    The 2048-token context window is managed by dropping oldest conversation turns in a band:
    when estimated history exceeds 90\% of context capacity, turns are evicted until it falls below 70\%,
    avoiding repeated KV cache invalidation on consecutive prompts.\footnote{%
      The combination of band-based eviction and llama.cpp's native context shifting is critical for interactive performance.
      Our initial implementation used the \texttt{llama-cpp-python} bindings,
      which in our implementation re-tokenized and re-evaluated the entire remaining history from scratch whenever the message list changed;
      on the Jetson Orin this added 30--60\,s of latency to the first prompt after every trim event.
      Switching to the \texttt{llama.cpp} HTTP server with its native KV cache shifting reduced post-trim TTFT to 1--3\,s.
      A naive single-threshold trim that fires on every subsequent prompt would compound this cost;
      the band ensures eviction happens at most once per several turns -- and is completed partially or entirely during VAD and STT.%
    }
    An Ollama~\cite{ollama} backend is also supported as an alternative runtime.
  \item \textbf{Content Filtering.}
    A streaming n-gram detector checks each token as it is generated,
    and each completed sentence is classified by Granite Guardian HAP-38M~\cite{padhi2024graniteguardian} running as a quantized INT8 ONNX graph on CPU.
    On match from either layer, generation halts immediately and the conversation context is reset.
    Both mechanisms are described in detail in Section~\ref{sec:safety}.
  \item \textbf{Text-to-Speech (TTS).}
    Piper TTS~\cite{piper_tts} synthesizes speech from text fragments,
    streamed to a persistent ALSA output to avoid underruns.
    A LaTeX-to-speech post-processor converts mathematical notation before synthesis.
\end{enumerate}

\noindent The pipeline is fully streaming:
TTS begins speaking the first sentence while the LLM is still generating subsequent ones.
Barge-in detection monitors microphone levels during playback,
allowing the child to interrupt at any sentence boundary.

\section{Hardware Platform}

Kian targets the NVIDIA Jetson Orin Nano (8\,GB unified memory).
The choice is driven by two properties of interactive voice systems:

\begin{itemize}[nosep]
  \item \textbf{TTFT over throughput.}
    In a conversational loop, the child waits from the moment they stop speaking until the assistant begins responding.
    What matters is how quickly the first token arrives, not the sustained generation rate (see Section~\ref{sec:latency}).
  \item \textbf{GPU with high memory bandwidth.}
    The Jetson Orin's unified memory architecture provides the bandwidth needed to keep TTFT low for quantized transformer inference,
    without discrete GPU memory transfer overhead.
\end{itemize}

\noindent Running headless (without a desktop environment) frees 1--1.3\,GB of RAM, improving GPU layer offload.
The total memory budget is approximately:

\begin{table}[H]
\centering
\begin{tabular}{lr}
\toprule
Component & Memory \\
\midrule
OS / system (headless) & $\sim$1.0\,GB \\
llama.cpp + LLM (Qwen3-4B Q4\_K\_M) & $\sim$2.5\,GB \\
Safety classifier (CPU) & $\sim$0.1\,GB \\
STT (Whisper base.en) & $\sim$0.2\,GB \\
VAD (Silero) + TTS (Piper) & $\sim$0.2\,GB \\
KV cache (2K context) & $\sim$0.3\,GB \\
\midrule
Headroom & $\sim$3.2\,GB \\
\bottomrule
\end{tabular}
\end{table}

\section{Speech-Only Control}

A core design principle is that no screen, keyboard, or companion device is required.
All configuration is performed through natural speech commands:

\begin{itemize}[nosep]
  \item \textbf{Identity:}
    ``Your name is Sophie'' or ``I'll call you Max'' changes the assistant's persona and persists the choice across sessions.
  \item \textbf{Grade level:}
    ``I'm in second grade'' or ``I'm in grade 5'' adjusts vocabulary and explanation complexity.
    Supports spoken numbers (``two'', ``second'') and digits alike, for grades pre-K through 12.
  \item \textbf{Volume:}
    ``Whisper'', ``inside voice'', ``speak up'', and ``shout'' map to four volume presets (30\%, 50\%, 80\%, 100\%) controlled via PulseAudio.
  \item \textbf{Context:}
    ``Let's start fresh'' or ``forget everything'' resets conversation history.
    Eight minutes of silence triggers an automatic reset.
  \item \textbf{Voice:}
    ``Change your voice'' or ``can I talk to someone else'' switches to a different TTS voice.
\end{itemize}

\noindent All settings---name, grade, voice, and volume---are persisted to disk and restored on startup.

\section{Adaptive Response Length}

A 4B-parameter model cannot reliably judge from a single instruction when to give a brief answer versus an elaborate explanation versus a vivid story.
To address this, the system prompt is rewritten before each LLM turn based on regex pattern matching of the user's input and whether Wikipedia reference material was retrieved:

\begin{itemize}[nosep]
  \item \textbf{Short} (default): 1--3 sentences.
    Triggered by simple questions and casual conversation (e.g., ``What is the quadratic formula?'').
  \item \textbf{Explain}: 50--150 words.
    Triggered by phrases such as ``explain'', ``teach me'', ``tell me about'', ``how does'',
    or by the presence of retrieved Wikipedia context.
    LaTeX math notation is permitted.
  \item \textbf{Story}: 100--300 words.
    Triggered by ``tell me a story'', ``describe a'', ``imagine a'', and similar creative prompts.
    LaTeX is suppressed to avoid injecting notation into narrative speech.
\end{itemize}

\noindent This approach moves the length decision out of the model and into deterministic pattern matching,
which is more reliable at this model scale.
The benchmarks in Section~\ref{sec:eval} use a fixed system prompt (short mode) to minimize output length disparities between models and different runs.

\section{Content Safety}\label{sec:safety}

Children's content filtering requires multiple layers, since no single mechanism is sufficient:

\begin{enumerate}[nosep]
  \item \textbf{System prompt.}
    The LLM is instructed to behave as an educational cartoon animal speaking to a child at a specific grade level.
    This biases generation toward age-appropriate content.
  \item \textbf{Streaming token filter.}
    A sliding-window n-gram matcher checks every token against a curated list of 100+ prohibited words and phrases
    spanning profanity, slurs, sexual content, violence, drug references, and adult/gambling websites.
    Detection is immediate---the response is discarded, the context is reset, and a neutral deflection is spoken.
  \item \textbf{Neural safety classifier.}
    Each complete sentence is classified by Granite Guardian HAP-38M~\cite{padhi2024graniteguardian},
    a 38.5M-parameter RoBERTa-based binary toxicity detector purpose-built for hate, abuse, and profanity detection.
    The model runs as a quantized INT8 ONNX graph on CPU, avoiding GPU memory contention with the primary LLM.
    Inference takes approximately 30\,ms per sentence.
    Unsafe sentences halt generation and reset the conversation context.
    A throttle mechanism ensures safety checks are interleaved with TTS playback rather than executed in bursts.
\end{enumerate}

\noindent The n-gram filter list is maintained as a ROT-13 obfuscated text file,
with helper scripts making it straightforward for educators and parents to review and extend.
The neural classifier provides contextual understanding that rule-based filters cannot---for example,
detecting abusive language composed entirely of individually benign words.

In practice, the n-gram filter reliably catches direct profanity and known slurs,
while the LLM classifier catches subtler cases such as violent scenarios described in clinical language.
The pipeline is weakest against indirect elicitation---a child who asks ``what words are bad?''
or requests a story with violent themes using age-appropriate vocabulary may receive content that neither filter flags.
Misspellings and phonetic substitutions (e.g., ``fck'') can also evade the n-gram matcher.
The system prompt biases the LLM away from these outputs, but a determined child can find gaps.
We consider this an inherent limitation of rule-based and small-model classifiers operating without human oversight.

\section{Mathematical Post-Processing}

Audio-based educational assistants face a fundamental challenge with mathematics:
spoken language is inherently sequential and ambiguous for mathematical expressions.
``a plus b over c'' could mean $a + \frac{b}{c}$ or $\frac{a+b}{c}$.

Kian addresses this by allowing the LLM to emit standard LaTeX notation,
which is converted to spoken English before TTS synthesis.
The system:

\begin{itemize}[nosep]
  \item Detects math regions via standard delimiters (\texttt{\$...\$}, \texttt{\$\$...\$\$}, \verb|\(...\)|, \verb|\[...\]|, and \verb|\begin{equation}|).
  \item Converts fractions, roots, superscripts, subscripts, Greek letters, operators, set theory notation, calculus operators, and quantifiers.
  \item Applies disambiguation rules:
    compound sub-expressions (those containing addition or subtraction) are wrapped with ``the quantity \ldots end quantity'' in contexts where boundaries are ambiguous (e.g., square roots),
    while fractions use ``over'' as a natural delimiter.
  \item Distinguishes unary negative from binary minus:
    $-x$ becomes ``negative x'' while $a - b$ becomes ``a minus b''.
  \item Separates adjacent single-letter variables (e.g., $ix$ becomes ``i x'') to prevent TTS from reading them as words,
    while preserving known multi-letter tokens like function names.
  \item Leaves prose outside math delimiters completely untouched,
    so punctuation like ``Hey there!'' is never misinterpreted as factorial.
\end{itemize}

\noindent This makes it feasible for a child to ask ``what is the quadratic formula?''
and hear a correct, unambiguous spoken rendering of $x = \frac{-b \pm \sqrt{b^2 - 4ac}}{2a}$.

\section{Latency Considerations}
\label{sec:latency}

In interactive voice applications, the relevant metric is \emph{turn latency}:
the time from when the child stops speaking to when the assistant begins responding.
This decomposes as:

$$t_{\text{turn}} = t_{\text{silence}} + t_{\text{STT}} + t_{\text{TTFT}} + t_{\text{TTS,first}}$$

\noindent Table~\ref{tab:latency} shows approximate values for each component on the target hardware using Qwen3-4B under llama.cpp.

\begin{table}[H]
\centering
\begin{tabular}{lr}
\toprule
Component & Time \\
\midrule
$t_{\text{silence}}$ (VAD silence window) & 1.2\,s \\
$t_{\text{STT}}$ (Whisper base.en transcription) & $\sim$0.3\,s \\
$t_{\text{TTFT}}$ (LLM first token, mean) & $\sim$0.17\,s \\
$t_{\text{TTS,first}}$ (Piper synthesis + playback start) & $\sim$0.1\,s \\
\midrule
Estimated turn latency & $\sim$1.8\,s \\
\bottomrule
\end{tabular}
\caption{End-to-end turn latency breakdown for Qwen3-4B (llama.cpp) on Jetson Orin Nano.
Values are engineering estimates from representative runs;
TTFT is the measured mean from Table~\ref{tab:perf}.}
\label{tab:latency}
\end{table}

\noindent The VAD silence window dominates turn latency;
all other components combined add under 0.6\,s.
Tokens per second is not the bottleneck---even at modest generation rates, text accumulates faster than it can be spoken.
The streaming architecture exploits this:
TTS begins on the first sentence boundary while the LLM continues generating.

\section{Wikipedia Context Augmentation}

Small language models hallucinate frequently on factual questions---a serious issue for an educational tool.
Kian mitigates this with a retrieval-augmented generation (RAG) pipeline backed by Simple English Wikipedia, an edition written at a reading level appropriate for children.

\subsection{Offline Index}

A build script (\texttt{build-wiki-db.py}) downloads the Simple Wikipedia XML dump (${\sim}350$\,MB)
and constructs a local SQLite database with an FTS5 full-text search index (${\sim}170$\,MB).
The entire index lives on-device; no network access is needed at query time.

\subsection{Topic Extraction}

Before each LLM turn, the child's utterance is matched against a set of 20+ regex patterns designed to extract educational topics from natural speech:
``what is a black hole'', ``tell me about volcanoes'', ``who was Marie Curie'', ``how does photosynthesis work'', and similar phrasings.
Stopwords are stripped, and the remaining terms are used as an FTS5 query.

\subsection{Retrieval and Injection}

If the FTS5 search returns a result above a BM25 relevance threshold,
the first 500 characters of the matching article are injected into the system prompt as reference material for that turn.
The LLM is instructed to use the reference only if relevant and never to mention its source.
When conversation turns are evicted from the context window,
corresponding wiki references are also evicted, keeping the context clean.

This approach grounds the LLM's responses in factual content without requiring a larger model or internet access.

\section{Model Evaluation}\label{sec:eval}

\subsection{Methodology}

We evaluated candidate models using a sequential prompt benchmark that simulates a child's conversation.
Twenty prompts progress from a simple greeting through factual questions to a creative storytelling request,
building conversational context incrementally and exceeding the 2048-token context window to exercise trim behavior.
Each conversation is scored on three dimensions by an LLM judge:

\begin{itemize}[nosep]
  \item \textbf{Overall}: holistic 1--5 conversation quality.
  \item \textbf{Fourth-wall breaks}: count of responses that broke the cartoon-animal persona
    (e.g., mentioning being an AI, or naming the persona type).
  \item \textbf{Factual errors}: count of responses on factual prompts (science, history, mythology) containing clear errors;
    imaginative prompts (stories, game suggestions, preferences) are excluded.
\end{itemize}

\noindent Models were benchmarked on the target hardware (Jetson Orin Nano, 8\,GB) using two runtimes:
llama.cpp~\cite{llamacpp} running as an OpenAI-compatible HTTP server, and
Ollama~\cite{ollama}, which also serves an OpenAI-compatible API.
Both runtimes use CUDA for GPU offload of quantized GGUF models.
The evaluated model families include
Qwen3~\cite{yang2025qwen3},
Llama~3.2~\cite{grattafiori2024llama},
Granite~3.3~\cite{granite2024} and 4.0~\cite{granite2025},
Ministral-3~\cite{liu2026ministral},
Nemotron-3 Nano~\cite{nemotron3nano},
and Bonsai-8B~\cite{bonsai8b} (a 1-bit quantized model).

\subsection{Performance Results}

Except where noted (the \texttt{-c1024} variant), all models were evaluated with a 2048-token context window across 5 benchmark runs (100 samples per model).
No temperature, top-$p$, or other sampling parameters were adjusted; all models used their runtime defaults.
Repeating the same prompt sequence multiple times averages over two sources of variance:
randomness in the model's own generated responses (due to sampling temperature),
and randomness in the LLM-as-judge evaluation of qualitative scores.
Table~\ref{tab:perf} summarizes latency and throughput,
with Table~\ref{tab:probe} covering short filler-word inputs;
Table~\ref{tab:quality} summarizes qualitative response quality.

\begin{table}[H]
\centering
\small
\begin{tabular}{l|r|r|r|r|r}
\toprule
Engine & Mean TTFT (s) & s.e.\ (s) & p95 TTFT (s) & tok/s & RAM (GB) \\
\midrule
ollama:Granite 3.3-2B & 0.228 & 0.017 & 0.346 & 25.9 & 1.8 \\
server:Bonsai-8B & 0.310 & 0.017 & 0.380 & 14.2 & 1.2 \\
ollama:Granite 4-3B & 0.354 & 0.014 & 0.447 & 18.4 & 2.4 \\
ollama:Qwen3-4B & 0.491 & 0.020 & 0.596 & 15.5 & 3.4 \\
ollama:Llama 3.2-3B & 0.527 & 0.020 & 0.622 & 19.1 & 2.5 \\
ollama:Ministral-3 3B & 0.570 & 0.017 & 0.658 & 19.5 & 4.8 \\
ollama:gemma3:4b & 0.885 & 0.020 & 1.010 & 15.6 & 4.3 \\
ollama:Qwen3.5-2B & 1.007 & 0.040 & 1.189 & 21.9 & 3.5 \\
ollama:Nemotron-3 Nano 4B & 1.019 & 0.108 & 1.544 & 15.7 & 5.2 \\
llamacpp:Granite 3.3-2B & 1.036 & 0.292 & 3.049 & 8.2 & 1.6 \\
llamacpp:Granite 4.0 H-Micro & 1.185 & 0.339 & 3.534 & 6.7 & 1.9 \\
llamacpp:Granite 4.0 Micro & 1.224 & 0.329 & 3.492 & 6.2 & 2.1 \\
llamacpp:Qwen3-4B & 1.468 & 0.403 & 4.236 & 5.1 & 2.5 \\
llamacpp:Granite 4.0 Micro IQ4 & 2.195 & 0.646 & 6.676 & 5.6 & 1.9 \\
llamacpp:Qwen3.5-2B & 4.899 & 0.777 & 8.241 & 8.3 & 1.3 \\
\bottomrule
\end{tabular}
\vspace{2pt}
\noindent{\footnotesize All models achieved 100\% GPU offload.\textsuperscript{$\dagger$}}
\noindent{\footnotesize \textsuperscript{$\dagger$}\,Kernel page cache from prior model loads (via \texttt{mmap}) was cleared between engines using \texttt{drop\_caches} to ensure full offload on Jetson unified memory.}

\caption{Performance benchmarks on Jetson Orin Nano (8\,GB, headless).
Steady-state and post-trim TTFT over 5 runs $\times$ 20 conversational prompts = 100 samples per engine.
Unless otherwise noted, models use Q4\_K\_M quantization and a 2048-token context.
All models achieved 100\% GPU offload;
the kernel page cache from prior model loads (via \texttt{mmap}) was cleared between engines with \texttt{drop\_caches} to ensure full offload on Jetson unified memory.}
\label{tab:perf}
\end{table}

\begin{table}[H]
\centering
\small
\begin{tabular}{l|r|r}
\toprule
Engine & Avg TTFT (s) & Max TTFT (s) \\
\midrule
server:ibm-granite\_granite-4.0-micro-IQ4\_XS & 0.165 & 0.796 \\
server:granite-4.0-micro-Q4\_K\_M & 0.233 & 1.428 \\
ollama:Granite 3.3-2B & 0.317 & 1.190 \\
server:granite-4.0-h-micro-Q4\_K\_M & 0.323 & 2.086 \\
server:qwen3-4b-instruct-2507-q4\_k\_m & 0.323 & 1.669 \\
server:granite-4.0-h-micro-Q4\_K\_M-c1024 & 0.349 & 1.218 \\
server:granite-3.3-2b-instruct-Q4\_K\_M & 0.495 & 1.314 \\
ollama:Granite 4-3B & 0.559 & 1.380 \\
server:Bonsai-8B & 0.570 & 2.899 \\
ollama:Qwen3-4B & 0.636 & 2.038 \\
ollama:Ministral-3 3B & 0.675 & 0.786 \\
ollama:Llama 3.2-3B & 0.685 & 1.425 \\
server:Qwen3.5-2B-Q4\_K\_M & 0.781 & 1.114 \\
ollama:gemma3:4b & 1.018 & 2.241 \\
ollama:Qwen3.5-2B & 1.306 & 1.761 \\
ollama:Nemotron-3 Nano 4B & 2.953 & 3.402 \\
\bottomrule
\end{tabular}

\caption{TTFT on short filler-word probe prompts (``Ummm'' and ``Okay''),
with responses capped at one token.
Each prompt arrives after 20 conversational turns,
so the history is already built up.
Probes that triggered a context trim are excluded so the figures isolate baseline inference latency on minimal user inputs.}
\label{tab:probe}
\end{table}

\begin{table}[h]
\centering
\small
\begin{tabular}{l|r|r|r|r}
\toprule
Engine & Err Avg & Err Max & Chatty & Poor Prose \\
\midrule
llamacpp:Qwen3-4B & 0.0 & 0 & 0.0 & 0.0 \\
ollama:Granite 3.3-2B & 0.0 & 0 & 2.0 & 2.0 \\
llamacpp:Granite 4.0 Micro IQ4 & 0.0 & 0 & 0.0 & 3.0 \\
llamacpp:Granite 4.0 H-Micro & 0.0 & 0 & 0.0 & 3.0 \\
ollama:Llama 3.2-3B & 1.0 & 1 & 0.0 & 0.0 \\
ollama:Ministral-3 3B & 1.0 & 1 & 3.0 & 0.0 \\
ollama:Granite 4-3B & 1.0 & 1 & 0.0 & 0.0 \\
ollama:gemma3:4b & 1.0 & 1 & 1.0 & 1.0 \\
server:Bonsai-8B & 1.0 & 1 & 0.0 & 1.0 \\
llamacpp:Qwen3.5-2B & 1.0 & 1 & 0.0 & 1.0 \\
ollama:Nemotron-3 Nano 4B & 1.0 & 1 & 0.0 & 2.0 \\
llamacpp:Granite 3.3-2B & 1.0 & 1 & 1.0 & 2.0 \\
llamacpp:Granite 4.0 Micro & 1.0 & 1 & 0.0 & 2.0 \\
ollama:Qwen3-4B & 3.0 & 3 & 1.0 & 1.0 \\
ollama:Qwen3.5-2B & 3.0 & 3 & 1.0 & 3.0 \\
\bottomrule
\end{tabular}
\caption{Qualitative response quality across 1 benchmark runs (8 prompts each). All metrics are lower-is-better: Err = factual errors per conversation, Chatty = verbose or repetitive responses (of 5), Poor Prose = weak explanatory/story responses (of 3). Evaluated by Claude Opus 4.6.}
\label{tab:quality}
\end{table}

\begin{table}[h]
\centering
\small
\begin{tabular}{l|l}
\toprule
Engine & Representative flub \\
\midrule
llamacpp:Qwen3-4B & --- \\
ollama:Granite 3.3-2B & --- \\
llamacpp:Granite 4.0 Micro IQ4 & --- \\
llamacpp:Granite 4.0 H-Micro & --- \\
ollama:Llama 3.2-3B & \textit{``about four times faster than a human''} \\
ollama:Ministral-3 3B & \textit{``That's even faster than a cheetah!''} \\
ollama:Granite 4-3B & \textit{``swims much faster than any running animal on land''} \\
ollama:gemma3:4b & \textit{``about 100 times faster than an average human''} \\
server:Bonsai-8B & \textit{``faster than a cheetah on land''} \\
llamacpp:Qwen3.5-2B & \textit{``can swim at speeds over 100 miles per hour''} \\
ollama:Nemotron-3 Nano 4B & \textit{``about 140 times faster than a human''} \\
llamacpp:Granite 3.3-2B & \textit{``the sailfish would likely win''} \\
llamacpp:Granite 4.0 Micro & \textit{``sailfish would likely win because it swims much faster''} \\
ollama:Qwen3-4B & \textit{``They also practice a lot''} \\
ollama:Qwen3.5-2B & \textit{``fastest fish is the tuna''} \\
\bottomrule
\end{tabular}
\caption{Representative factual errors from qualitative evaluation.}
\label{tab:flubs}
\end{table}

\antml:thinking>
The user wants a short discussion section for a LaTeX litepaper, analyzing the qualitative evaluation table. Let me write 1-2 concise paragraphs highlighting Granite, Qwen3, and Llama families, noting trade-offs.
</thinking>

The qualitative evaluation reveals that factual accuracy and prose quality are largely independent axes, and no single model family dominates both. \textbf{Qwen3-4B} (llama.cpp) is the only model to achieve a perfect score across all four metrics---zero errors, zero chattiness, and no poor prose---making it the strongest overall candidate at this scale. However, its smaller sibling \textbf{Qwen3.5-2B} is notably inconsistent: the llama.cpp variant produces only one error but suffers from weak prose (PoorProse = 1), while the Ollama variant degrades sharply to three errors and three poor-prose marks, suggesting that quantization and runtime choice materially affect output quality for smaller models. The \textbf{IBM Granite} family presents a different trade-off. Granite 3.3-2B and the Granite 4.0 Micro variants achieve zero factual errors, but all three receive PoorProse scores of 2--3, indicating that their longer-form responses lack the fluency expected of a children's voice assistant. \textbf{Granite 4-3B} improves prose (PoorProse = 0) but introduces one factual error. \textbf{Meta Llama 3.2-3B} strikes an appealing balance---one error, zero chattiness, and zero poor prose---though its single flub (``about four times faster than a human'') is relatively mild compared to the more egregious claims produced by other models.

These results suggest two viable strategies for deployment on the Jetson Orin Nano. For applications where factual reliability is paramount and slightly stiff prose is acceptable, the Granite 4.0 Micro family (IQ4\_XS or half-precision) offers zero-error performance at just 2B parameters, the smallest footprint tested. Where conversational naturalness matters more---as it likely does for young children---Qwen3-4B and Llama 3.2-3B are preferable, accepting a small error risk in exchange for markedly better prose. Notably, the \textbf{Ministral-3B} model's high chattiness score (3.0) and \textbf{Gemma 3 4B}'s mixed results (Chatty = 1, PoorProse = 1) make them less competitive despite comparable parameter counts. A single evaluation conversation per model limits statistical power, but the consistency of these patterns across model families lends confidence to the ranking.


\section{Limitations and Future Work}

\begin{itemize}[nosep]
  \item \textbf{Language support} is currently English-only across all pipeline stages.
  \item \textbf{Content filtering} combines rule-based n-gram matching with a neural toxicity classifier,
    but neither mechanism can guarantee detection of novel circumventions or subtly inappropriate content.
    Adversarial probing by children is a realistic threat model.
  \item \textbf{Mathematical rendering} handles standard LaTeX but not all edge cases
    (e.g., matrices, piecewise functions, or deeply nested expressions).
  \item \textbf{No visual output} means that some mathematical concepts (geometry, graphs) are difficult to convey through speech alone.
\end{itemize}

\section{Conclusion}

Kian demonstrates that a practical, fully offline educational voice assistant can run on modest edge hardware.
Recent advances in small-model quality and quantization make this class of system newly feasible:
a 4-bit 4B-parameter model on an 8\,GB Jetson achieves sub-200\,ms mean time-to-first-token with no factual errors over our internal benchmark run,
while leaving enough headroom for a full speech pipeline and safety classifier.

The architectural bet is that privacy, connectivity independence, and child autonomy are worth the tradeoffs inherent in small on-device models---%
limited context windows, narrower world knowledge, and imperfect content filtering.
For an educational voice assistant aimed at children, we believe those tradeoffs favor the on-device approach:
the privacy guarantees are absolute rather than policy-dependent,
the system works anywhere a power outlet exists,
and a child can use it without adult mediation.
The fully offline approach is likely to remain relevant regardless of how cloud-based assistants evolve,
as its privacy and accessibility advantages are structural, not just a current-generation artifact.

\section*{Acknowledgments}

This paper and the accompanying codebase were composed with assistance from Claude Opus 4.6 (Anthropic).

\bibliographystyle{unsrt}
\bibliography{refs}

\end{document}
